\documentclass[11pt]{article}
\usepackage[margin=25mm]{geometry}
\usepackage{amsmath,amssymb,hyperref,booktabs}
\usepackage[T1]{fontenc}
\title{Symbolic synthesis: global minimum certificates and an objective obstruction}
\author{Computational proof ledger}
\date{7 November 2026}
\begin{document}\maketitle
\noindent The certificates below distinguish exact circuit identities, global
entangler minima, and global minima of the parameter-sensitive objective.
These are separate statements. They do not establish optimality throughout
all $327\times32$ target/basis cells.

\section{An obstruction in the unrestricted historical objective}
For an operation with $k$ wires and parameters $p_j$, the historical score is
\[
 2^k+\frac15\sum_j\left(\log_2|\operatorname{den}(p_j)|+\frac1{100}\right),
\]
where numerator and denominator are taken after symbolic simplification.
This expression need not be nonnegative for unrestricted real expressions.
Let $d=\log(2049/2048)$ and $a=1/d$. The native one-qubit word
\[
 U_3(-a,0,0)U_3(a,0,0)=I
\]
is an exact matrix identity, not merely equality up to scalar. Its score is
\[
 C=\frac{1003}{250}+\frac25\log_2 d < -\frac{97}{250}<0.
\]
Indeed $0<\log(1+x)<x$ for $x>0$, since
$[\log(1+x)-x]'=-x/(1+x)<0$, and here $x=2^{-11}$.
Appending arbitrarily many copies of the word preserves the target, phase,
and entangler count, while its score tends to $-\infty$. Consequently there
is \emph{no unrestricted global minimum} of the secondary historical score
in bases admitting these U3 words. No amount of additional search proves a
minimum that does not exist.

One concrete repair is to replace the denominator term by
$\log_2\max(1,|\operatorname{den}(p_j)|)$. An exact compatibility audit of the
1,339 distinct retained parameter expressions in the pre-refinement archive
found no changed scores and no undecided denominator comparisons.
The user approved this repair and the compiler now uses objective version 2.
The expanded audit covered 1,390 distinct retained expressions, with zero changed or undecided scores. The archive was migrated after this check; historical source snapshots retain their original contracts.
Alternatively the admissible expression language can be explicitly restricted;
any resulting minimum claim must name that restriction.
Under version 2 every local costs at least two and every entangler at least four.
An incumbent score $C$ at fixed entangler count $E$ therefore bounds every potentially cheaper word to at most $\lfloor(C-4E)/2\rfloor$ locals. This bounds word length; continuous angle optimality still needs a separate proof.
The historical obstruction does not affect entangler-count lower bounds or the finite
alphabet proof in the next section.

\section{Global phase-resolved Clifford minima in Clifford+T}
The supplied alphabet is exactly $\{H,S,T,T^\dagger\}$. Each fixed gate has
historical score two, so minimum score is equivalent to minimum word length.
Every matrix is stored exactly in $\mathbb Z[\sqrt2,i]/2^k$, with common powers
of two removed from its integer representation. Thus hashing and equality
use integer arithmetic; floating comparisons play no role.

Breadth-first search through radius seven reaches 628 distinct matrices,
including non-Clifford intermediate states. For each of the 192 exact-phase
Clifford targets $G$, enumerate each first half $A$ in that ball and look up
$GA^{-1}$ in the same ball. The resulting witnesses have maximum length 14.
Every shorter word can be split into two halves of length at most seven;
therefore the lookup excludes every cheaper word in the supplied alphabet.
The proof covers all 192 literal phases, and taking the minimum over the eight
phases of each representative proves the corresponding projective minima.
It extends the H/S proof without assuming that intermediate gates are Clifford.

For projective completeness, every H/S/T/Tdg matrix has entries in
$K=\mathbb Q(\sqrt2,i)=\mathbb Q(\zeta_8)$, of degree four. The determinant
of any word is an eighth root; the determinant of a Clifford target is a
fourth root. If a word is $\lambda C$, then $\lambda\in K$ and
$\lambda^2$ is an eighth root. An odd sixteenth root has degree eight,
so cannot lie in $K$. Thus $\lambda$ is one of the eight phases already
enumerated. This excludes additional scalar phases from the projective
search, including words with non-Clifford intermediates.


The canonical ball SHA256 is
\begin{quote}\small\ttfamily
 e7c7e0492f5ee55d0b80adacc65331c9dc0c2230cd04e2b1b53c58441e7462af.
\end{quote}
The implementation replays all target witnesses and checks their distances
against the Clifford shortest-path certificate. Altered proof radii,
distances, alphabets, and missing global certificates are rejected.
Continuous alphabets remain outside this finite-word proof.

\section{Two B gates for the full Cartan family}
Use the convention
\[
 A(x,y,z)=\exp\{i(xXX+yYY+zZZ)\},\qquad B=A(\pi/4,\pi/8,0).
\]
The two-B upper bound is established in the literature~\cite{b}. The new
implementation supplies explicit symbolic outer frames, all real-angle
branches, native U3 lowering, and physical scalar absorption.
Fold $y,z$ into $[-\pi/4,\pi/4)$ using $\pi/2$ shifts. Put
\[
 u=\sqrt2\sin y\cos z,\quad v=\sqrt2\cos y\sin z,\quad
 w=\sqrt{\cos2y\cos2z}.
\]
Then $u^2+v^2+w^2=1$. On the regular chart choose
\[
 2p_- =\operatorname{atan2}(-w,\sin(y-z)),\quad
 2p_+ =\operatorname{atan2}(-w,\sin(y+z)),
\]
\[
 \beta=\operatorname{atan2}(w,u),\qquad
 a=2\operatorname{atan2}(v,\sqrt{u^2+w^2}),\quad
 \chi=-p_--p_+-\pi.
\]
The circuit in time order has local layers
\[
\begin{array}{c|cc}
 &q_0&q_1\\\hline
\text{before first B}&U_3(\pi/2,\chi,0)&U_3(\pi+p_--p_+,-\pi/2,-\pi/2)\\
\text{between B gates}&U_3(-2x,0,0)&U_3(-a,-\beta,-\beta)\\
\text{after second B}&U_3(-\pi/2,0,\chi-\pi)&U_3(\pi-p_-+p_+,-\pi/2,\pi/2)
\end{array}
\]
and scalar $\exp\{i(\beta+p_-+p_+)\}$ on the folded chart.
An extra native U3 absorbs this scalar for literal equality. No GP primitive
is introduced. Local Y/Z frames restore the folds; the implementation also
absorbs the input/output Z gauges of PhasedFSim and its scalar.

The full matrix identity reduces six nonzero Pauli coefficients modulo
$c_y^2+s_y^2=1$, $c_z^2+s_z^2=1$ and
$w^2=(c_y^2-s_y^2)(c_z^2-s_z^2)$. The four chart corners and the zero-radius
branch are treated explicitly. Arbitrary $x$ follows from exact Pauli transport
through B. Parameter substitutions are regression checks, not family proofs.

For the lower bound, the phase-normalized spin-flip invariants~\cite{makhlin}
are
\begin{align*}
 I_1&=16(\cos2x\cos2y\cos2z+i\sin2x\sin2y\sin2z)^2,\\
 I_2&=4(\cos4x\cos4y\cos4z+i\sin4x\sin4y\sin4z).
\end{align*}
Local products have $I_1=16$; one B pulse has $I_1=I_2=0$.
Consequently the construction is globally minimum in entangler count whenever
\[
 I_1\ne16\quad\text{and}\quad(I_1\ne0\ \text{or}\ I_2\ne0).
\]
The certificate records this condition rather than claiming every specialization
needs two pulses. Scalar relaxation in the lower bound makes it valid for
both phase policies. No minimum local score is claimed for these charts.


\section{Additional analytic score floors and local transfers}
At zero entanglers, each wire on which the target acts nontrivially requires
at least one one-wire gate. An exact nonzero commutator with an embedded X
or Z proves such support. For a family, a nonzero exact specialization is a
witness that the entire family cannot omit that wire; cheaper exceptional
specializations remain possible. Literal nonidentity scalar targets also need
at least one physical operation. Combining these support bounds with the
nonnegative primitive floors certifies a secondary minimum whenever the
constructed score meets the bound. Optional default parameters are excluded
from the minimum mandatory-arity surcharge.

The current archive records 372 literal and 1,144 projective global
lexicographic minimum certificates from this floor argument. There are also
24 literal and 31 projective matches of the secondary floor at a fixed
entangler count whose entangler minimum is not established by this argument.
These counts can overlap existing Clifford and entangler certificates.

The C4X native-CZ transfer first proves the complete 36-CZ, 46-U3 parent,
then proves every distinct local replacement as a two-by-two identity.
The free-angle identity that Rz commutes with CZ on either wire justifies
collecting diagonal phases across interactions. This yields 69 locals in
the CZ/Rz/Rx basis, compared with the stored SDK's 72, and 130 locals in the
IBM CZ basis, compared with 143. The IBM construction searches exact native
Euler variants while carrying diagonal phases through CZ gates; each variant
and every removed full turn have their own exact identity.
Both transfers preserve 36 entanglers and a recorded projective scalar;
these are constructive improvements, not minimum-count proofs for C4X.


\section{Exact local polar fusion}
For a constant one-qubit unitary $M=\left(\begin{smallmatrix}A&B\\C&D\end{smallmatrix}\right)$,
put $r=|A|$, $s=|C|$. If $A,C\ne0$, define
\[
 \delta=\arg A,\quad \vartheta=2\operatorname{atan2}(s,r),\quad
 \varphi=\arg C-\delta,\quad \lambda=\arg(-B)-\delta.
\]
Then $M=e^{i\delta}U_3(\vartheta,\varphi,\lambda)$. Unitarity gives
$r^2+s^2=1$, $|B|=s$, and $\bar A B+\bar C D=0$.
The defining polar identities reproduce the first three entries; the last
is fixed by column orthogonality. If $A=0$, use
$\vartheta=\pi$, $\varphi=\arg C$, $\lambda=\arg(-B)$ and $\delta=0$.
If $C=0$, use $\vartheta=\varphi=0$ and
$\lambda=\arg D-\arg A$. Scalar blocks need no physical local operation.
These explicit charts avoid undefined arguments at zero.

The implementation proves primitive unitarity for arbitrary real arguments,
then uses closure under multiplication for every source block. Exact zero
predicates and algebraic minimal polynomials select the constant charts.
The complete parent circuit is proved separately. The local proof therefore
does not rely on numerical fidelities or a general simplifier rediscovering
polar branch identities. Rz/Ry lowering and literal U3 phase absorption have
independent free-angle matrix identities.

This reduces C3SX/iSWAP to 40 locals projectively and 41 literally, compared
with the SDK's projective 56. C4X/iSWAP uses 78 and 79, compared with SDK 113.
Their entangler counts remain 26 and 54. Peres/Rxx/Rz/Ry now matches the
SDK's 35 locals at five entanglers. These are upper bounds, not unrestricted
local-minimum claims. Only C3X/Rzz/U3 retains a higher local count in the
matched SDK audit: 28 locals score 115.594, whereas its available proved
24-local alternative scores 116.37. This establishes a tradeoff between
those two witnesses, not the impossibility of a better unknown circuit.

\section{Independent native lower bounds}
For a two-qubit unitary $U$, write $\Gamma=U(YY)U^T(YY)$ and $\Delta=\det U$.
The scalar-normalized local invariants used here are~\cite{makhlin}
\[
 g_1=\frac{(\operatorname{tr}\Gamma)^2}{16\Delta},\qquad
 g_2=\frac{(\operatorname{tr}\Gamma)^2-\operatorname{tr}(\Gamma^2)}{4\Delta}.
\]
They require no choice of a fourth root of $\det U$. An exact nonzero
realignment minor excludes a local product. Unequal invariants or unequal
fixed operator-Schmidt ranks exclude one fixed native pulse. For a variable
Ising interaction, every pulse has crossing Schmidt rank at most two;
a target of rank four therefore needs at least two pulses.

In the magic basis, arbitrary determinant-one local frames are real matrices
$O\in SO(4)$. Partition $O$ into four $2\times2$ blocks. Their squared
Frobenius norms are $X,Y,Y,X$, by subtracting the block row and column norm
equations. For two Ising pulses, with possibly different interaction angles,
\[
 \operatorname{tr}(D_1^2 O^T D_2^2 O)
 =2\cos(a+b)X+2\cos(a-b)Y\in\mathbb R.
\]
For the iSWAP edge the second sign becomes negative and the trace is again
real. Thus two such pulses require $g_1\geq0$ and real. Two square-root-SWAP
pulses instead have purely imaginary normalized spin trace, requiring
$g_1\leq0$ and real. The native zero- and one-pulse classes are excluded
separately before a three-pulse lower bound is asserted. This distinction is
necessary: a single SYC or square-root-SWAP need not be cancellable by a
second identical native pulse.

A single fSim has a monic cubic annihilator for its spin matrix,
\[
 (\Gamma-e^{-i\phi}I)(\Gamma^2-2\cos(2\theta)\Gamma+I)=0.
\]
The identity is checked for free real parameters, including degeneracies.
A target with four linearly independent vectorized powers
$I,\Gamma,\Gamma^2,\Gamma^3$ cannot be locally equivalent to one fSim.
One exact nonzero $4\times4$ Krylov minor certifies that independence.

For square-root-SWAP there is a stronger two-pulse relation. Its normalized
spin spectrum is a scalar multiple of a real reflection. A product of two
such reflections gives $g_1=-\rho^2$, $g_2=1-4\rho$, $0\leq\rho\leq1$.
Consequently
\[
 16g_1+(g_2-1)^2=0
\]
is necessary for two pulses. The implementation proves the reflection traces
by polynomial reduction modulo the unit-vector norm equation.

\section{An algebraic SQiSW two-pulse exclusion}
Let $V$ be a two-SQiSW circuit in the magic basis, with arbitrary real
quaternion middle factors
$aI+i(bX+cY+dZ)$ and $eI+i(fX+hY+kZ)$, each of unit norm.
Put $A=\det\operatorname{Re}V$, $C=\det\operatorname{Im}V$. Direct exact
polynomial identities give
\[
 C=\tfrac14\big[(a^2+d^2)(f^2+h^2)-(b^2+c^2)(e^2+k^2)\big]^2\geq0,
\]
\[
 \tfrac14\operatorname{Re}\operatorname{tr}(V^TV)
 =(ak-de)^2+(bh-cf)^2\geq0.
\]
For an SU(4) unitary, the latter trace is $A-C$ and
$A+C=(g_2+1)/4$; these relations are proved on a free Cartan diagonal and
transported through the real determinant-one KAK frames. Hence $A,C\geq0$.
Writing $g_1=(u+iv)^2$ gives
$u^2=(|g_1|+\operatorname{Re}g_1)/2$. Therefore two SQiSW pulses necessarily obey
\[
 W=g_2+1\geq0,\qquad W^2\geq8(|g_1|+\operatorname{Re}g_1).
\]
These are scalar-invariant tests; multiplication by $i$ merely interchanges
$A$ and $C$. The local and one-SQiSW classes also obey both inequalities.
An exact negative $W$ or negative squared gap therefore proves a three-pulse
lower bound directly, without numerical Weyl canonicalization. This section
does not supply a missing three-pulse construction; the uniform upper bound
comes from the published SQiSW compilation scheme~\cite{sqisw}.

For a parameterized target, one exact hard specialization proves a necessary
\emph{uniform family budget}. A matching circuit for the entire family proves
that budget optimal. Other parameter values can have smaller minima.
The archive stores this scope separately from fixed-target minima and leaves
the original construction operations, scores, and source indices unchanged.
No minimum local score follows from these native-count certificates.

The independent native audit matches 3,274 archived phase modes: 681 literal
and 928 projective fixed-target minima, and 710 literal and 955 projective
minimum uniform family budgets. Of these, 1,224 previously had neither a
minimum-count flag nor a conditional minimum-count certificate. The separate
exchange audit supplies 18 matched fixed-target certificates and 132 further
lower bounds without matching constructions; these counts can overlap the
native audit. The reports retain unsuccessful proof attempts separately.

\section{Reproduction}
Run \nolinkurl{pytest tests/symbolic/test_b_cartan.py}
and \nolinkurl{pytest tests/symbolic/test_clifford_t_global.py}.
The counterexample is regenerated by
\nolinkurl{symbolic_cost_proof.negative_identity_certificate()}.
Raw source-pinned archive refinements retain their exact implementation hashes
and separate phase policies. A verified circuit is an upper bound until its
matching lower bound and assumptions are recorded.
The independent audit is generated, replayed, and attached with
\nolinkurl{examples/decomposition/symbolic_entangler_audit.py}.
Its source-pinned report records successful certificates, declined cases,
and proof timeouts separately. These timeouts concern new lower-bound
searches; they do not undo the existing exact construction proofs.
\nolinkurl{examples/decomposition/symbolic_exchange_audit.py} adds the
separately replayed exchange witnesses. The original 790 construction and
Clifford certificate replays remain in \nolinkurl{verification.json}.

\begin{thebibliography}{9}
\bibitem{b} J. Zhang, J. Vala, S. Sastry, K. B. Whaley,
Minimum construction of two-qubit quantum operations,
Phys. Rev. Lett. 93, 020502 (2004),
\url{https://arxiv.org/abs/quant-ph/0312193}.
\bibitem{makhlin} Y. Makhlin, Nonlocal properties of two-qubit gates and mixed
states and optimization of quantum computations,
Quantum Information Processing 1, 243--252 (2002),
\url{https://arxiv.org/abs/quant-ph/0002045}.
\bibitem{sqisw} C. Huang et al., Quantum Instruction Set Design for Performance,
Phys. Rev. Lett. 130, 070601 (2023),
\url{https://arxiv.org/abs/2105.06074}.
\end{thebibliography}
\end{document}
